\chapter{The topological classification and algebraic witnesses}
\label{tx:sec:model-theory}
This chapter expands Section~11 of the paper. The main description of
T1--T4 is the topological classification in
Section~\ref{tx:sec:topological-spectrum}: empty, full, nonempty proper
clopen, or closed but not open hardness spectrum. Its proof uses eventual
prime-power periodicity and compactness, and applies to all the uniform
requests covered by the theorem.
We then distinguish algebraic witnesses from actual models. The final
Euler-structure discussion is optional historical background; it is
not needed for this classification and supplies no converse criterion.
A uniform target invites a semantic question: does every solution of
$\Gamma_n$ satisfy $G^{(n)}=0$? Our theorem classifies syntactic algebraic
derivation, not that semantic question. The two are related through
linear witnesses, but the relation is not an identification of the four
types with four classes of ordinary models.
Keep a fixed uniform axiom family $\Gamma_n$ and its fixed target request
over a field $K$ of prime characteristic $p$. Its profile $h_S(n)$ is
the maximum least degree needed for those targets in NS or ordinary PC.

\section{The four types as shapes of a hardness spectrum}
\label{tx:sec:topological-spectrum}
The current paper gives a useful picture of the classification. This
picture applies to \emph{every} aggregate or fixed-orbit uniform request
in the theorem, over an arbitrary field of characteristic $p$. It is
not restricted to Boolean equations or to a particular translation of
first-order logic. Its input is eventual prime-power periodicity of
bounded-degree derivability.

The residue-class structure is already implicit in Riis's 2009
report~\cite[Section~4, pp.~18--19]{RiisLong}. The following reformulates
that structure and extends its reading to the uniform consequences
considered here; it is not a claim to have newly discovered the four
cases. No converse to the Euler-model criterion is needed.

\subsection{A \texorpdfstring{$p$}{p}-adic point is a compatible address in a residue tree}
At level one record a residue modulo $p$; at level two record one modulo
$p^2$ which reduces to the first residue; continue in this way. A
compatible infinite sequence is a $p$-adic integer. For example, in
characteristic two, the successive addresses of the ordinary integer
$5$ are $1\bmod2$, $1\bmod4$, $5\bmod8$, $5\bmod16$, and so on.
Not every infinite address represents a nonnegative ordinary integer.
This distinction is precisely what makes the extra points useful.
For further background see Gouv\^ea~\cite{Gouvea}; the argument
needed here is supplied below.

A neighbourhood specifies a finite part of the address:
$a+p^L\mathbb Z_p$ contains all points with the chosen residue modulo
$p^L$. Such a neighbourhood is both open and closed, or \emph{clopen}:
its complement is the finite union of the other residue classes.
The whole space $\mathbb Z_p$ is compact. One way to see why is to
view its points as infinite paths in the finitely branching residue
tree. If a cover by residue neighbourhoods had no finite subcover,
one could choose a residue at each level whose descendants still had
no finite subcover. The resulting infinite path would be uncovered,
a contradiction. General open covers reduce to these basic ones.

\subsection{Colour the residues which are eventually easy}
Use exactly $N(d)$ and $Q(d)$ from \eqref{tx:eq:NQ}, including the
separate convention $N(0)=n_*$ and $Q(0)=1$. Define
\[
 A_d(h)=\{a\in\{0,\ldots,Q(d)-1\}:\ h(n)\le d
       \text{ for all }n\ge N(d),\ n\equiv a\pmod{Q(d)}\},
\]
\[
 \mathcal E_d(h)=\bigcup_{a\in A_d(h)}(a+Q(d)\mathbb Z_p),\qquad
 \mathcal H_p(h)=\mathbb Z_p\setminus\bigcup_{d\ge0}\mathcal E_d(h).
\]
Here $A_d(h)$ is a \emph{set of residues}, whereas $N(d)$ is a
\emph{size threshold}. The calligraphic $\mathcal E_d(h)$ is a subset
of $\mathbb Z_p$; it is not the set $E(h)$ of recurring degree values.
The sets $\mathcal E_d(h)$ increase with $d$. To check this, pass to
the common finer period and compare sufficiently late positive integer
representatives. Each set is clopen, so the hardness spectrum
$\mathcal H_p(h)$ is closed.

A point lies \emph{outside} this spectrum precisely when it has a
residue neighbourhood on which all sufficiently large integer sizes
have one uniform finite degree bound. This equivalent description
shows that choosing later valid onsets or finer valid prime-power
periods does not change the spectrum. Finitely many exceptional sizes
also make no difference. We have not defined a polynomial system at
a noninteger $p$-adic size: the point records compatible information
about ordinary finite sizes.

\begin{proposition}[The four shapes]\label{tx:prop:topological-spectrum}
For the profiles in the four-way theorem,
\begin{center}
\begin{tabular}{@{}ll@{}}
\toprule
Type & Shape of $\mathcal H_p(h)$\\
\midrule
$\mathrm T_1$ & Empty\\
$\mathrm T_2$ & All of $\mathbb Z_p$\\
$\mathrm T_3$ & Nonempty, proper and clopen\\
$\mathrm T_4$ & Closed but not open\\
\bottomrule
\end{tabular}
\end{center}
For the same uniform request, NS and ordinary PC have the same
hardness spectrum, even if their exact degree profiles differ.
\end{proposition}
\begin{proof}
If the spectrum is empty, the increasing open sets $\mathcal E_d$
cover compact $\mathbb Z_p$. A finite subcover is contained in one
$\mathcal E_D$, so that set is the whole space. This gives an eventual
degree bound $D$, exactly T1. Conversely such a bound makes the
spectrum empty. The spectrum is the whole space exactly when every
$\mathcal E_d$ is empty. This means that no finite degree recurs,
exactly T2.

In T3 the classification theorem gives a fixed nonempty proper union
of easy residue classes, with a uniform bound $D$. No bounded-degree
residue neighbourhood can intersect the linearly hard region: a common
refinement would contain arbitrarily large integers satisfying both
incompatible bounds. Thus $\bigcup_d\mathcal E_d=\mathcal E_D$,
and the spectrum is nonempty, proper and clopen.
Conversely, for a nonempty proper clopen spectrum its complement is
compact. The same finite-subcover argument makes the complement equal
to some $\mathcal E_D$. Any recurring finite value $e$ occupies a
residue neighbourhood on a tail, by truncated periodicity. That
neighbourhood belongs to $\mathcal E_e\subseteq\mathcal E_D$, whence
$e\le D$. At least one finite value recurs, and the nonempty spectrum
excludes T1. This is T3. The remaining closed-but-not-open possibility
is therefore exactly T4.

Finally $h_{\PC}(n)\le h_{\NS}(n)\le F(h_{\PC}(n))$, with $F$
independent of $n$. A residue neighbourhood is eventually bounded for
one system if and only if it is eventually bounded for the other.
This proves equality of their spectra; it does not prove equality of
$\mathcal E_d$ for a fixed $d$.
\end{proof}

\subsection{A worked binary example: hard at the limit, finite at every size}
Take the abstract profile $h(n)=2^{v_2(n)}$ for positive integers $n$.
Thus odd sizes have value one, sizes twice an odd number have value
two, and sizes four times an odd number have value four. This is the
toy profile already used in Chapter~\ref{tx:ch:profiles}, not an
assertion that a particular polynomial family has exactly these degrees.
The easy regions are
\begin{center}
\begin{tabular}{@{}rll@{}}
\toprule
Budget $d$ & $\mathcal E_d(h)$ & Still outside this easy region\\
\midrule
$0$ & $\varnothing$ & $\mathbb Z_2$\\
$1$ & $1+2\mathbb Z_2$ & $2\mathbb Z_2$\\
$2,3$ & $\mathbb Z_2\setminus4\mathbb Z_2$ & $4\mathbb Z_2$\\
$4,\ldots,7$ & $\mathbb Z_2\setminus8\mathbb Z_2$ & $8\mathbb Z_2$\\
$8,\ldots,15$ & $\mathbb Z_2\setminus16\mathbb Z_2$ & $16\mathbb Z_2$\\
\bottomrule
\end{tabular}
\end{center}
Consequently $\mathcal H_2(h)=\bigcap_L2^L\mathbb Z_2=\{0\}$.
Every positive size has finite value, yet there is no bound valid on
any neighbourhood of this limiting address. For every $L$, the
neighbourhood $2^L\mathbb Z_2$ contains the arbitrarily large positive
sizes $2^j$ with values $2^j$. It also contains the arbitrarily large
sizes $2^L(2k+1)$ of the fixed finite value $2^L$.
Thus each neighbourhood of zero meets the easy region but is not
uniformly easy. The singleton $\{0\}$ is closed and not open.
This is the T4 picture. The point $0\in\mathbb Z_2$ is an address,
not a new instance with an empty underlying universe.

For comparison, the abstract profiles $h(n)=1$ and $h(n)=n$ give
empty and full spectra. The profile which is $1$ on odd sizes and
$n$ on even sizes gives $2\mathbb Z_2$, the T3 shape. These examples
separate the four pictures without claiming a realisation theorem
for arbitrary closed subsets of $\mathbb Z_p$.

\subsection{What topology explains, and what still needs the quantitative proof}
In T3, looking at a single finite number of residue digits suffices
to separate the easy and hard regions. In T4 it never suffices near
a boundary point. Compactness is the step turning local finite bounds
into one global bound in T1, or one bound on the entire easy part in T3.
It does not compute that bound.

Topology alone does not give linear lower bounds. For example,
$g(n)=1+v_2(n)$ has the same singleton spectrum as the preceding toy
profile, but $g(n)\le1+\log_2 n$. Its threshold periods grow
exponentially with the degree budget, rather than obeying the linear
period bound proved here. Thus it is not a counterexample to the
paper's theorem. The estimates $N(d)=O(d)$ and $Q(d)=O(d)$ are what
produce the linear gap and finite linear peaks. Keep the qualitative
picture and that quantitative input distinct.

\section{A quotient remembers what remains unproved}
For a fixed degree budget $d$, the ambient polynomial space must also
contain the target. We therefore allow polynomials up to degree
$\max\{d,D_G\}$, where $D_G$ is the fixed target-degree bound.
Form its quotient by $I_d(n)$ for NS or by $C_d(n)$ for PC, the consequence spaces
of Proposition~\ref{tx:prop:closure}. A target is derivable exactly when
its quotient class is zero. If the class is nonzero, elementary linear
algebra separates it from zero by a functional. This is a completely
abstract description of the failure to prove a consequence.

For target one the separating functional is normalised by $\lambda(1)=1$.
For a general target it must instead satisfy $\lambda(G)=1$.
That small change matters: a witness for failure to derive $x$ need not
be normalised on one at all. Nor is it generally evaluation at a model.
The two examples below make the distinction explicit.

\section{Dual witnesses for failure to derive a target}
For refutations, the design viewpoint of Buss~\cite{BussDesign}
organises such linear witnesses. We state the elementary separation
criterion for a prescribed target, whose normalisation is on that
target rather than necessarily on the constant $1$.
For either system write $W_d(n)$ for its degree-$d$ consequence space:
$W_d(n)=C_d(n)$ for PC and $W_d(n)=I_d(n)$ for NS. Let
$\delta=D_G$ be the fixed upper bound on target degree. We view
$W_d(n)$ inside $P_{\le\max\{d,\delta\}}(n)$, so the ambient space
also contains every requested target even when its degree exceeds $d$.

The \emph{dual} of a vector space is its space of linear functionals.
A functional \emph{annihilates} a subspace when it is zero on every
vector of that subspace. Thus the two conditions below say that the
functional ignores all proved consequences but detects the chosen
target. The condition on the target fixes the scalar normalisation of the
functional; it does not make the functional multiplicative.

\begin{definition}[A degree-bounded target witness]
\label{tx:def:consistency-witness}
For a requested target $G_b^{(n)}$, a degree-$d$ witness is a $K$-linear
functional
\[
 \lambda:P_{\le\max\{d,D_G\}}(n)\longrightarrow K,\qquad
 \lambda(G_b^{(n)})=1,\qquad\lambda(W_d(n))=0.
\]
For an all-target request, a witness consists of an index $b$ and such
a functional. For a single fixed orbit, one may use any available
representative. With target $1$, this reduces to the usual normalisation
$\lambda(1)=1$ for a bounded-degree non-refutability witness.
\end{definition}

\begin{proposition}[Dual form of nonderivability]
\label{tx:prop:dual-witness}
A degree-$d$ target witness exists exactly when the target has no
derivation of degree at most $d$ in the chosen system. For an all-target
request it exists exactly when $h(n)>d$.
\end{proposition}
\begin{proof}
If $G_b^{(n)}\in W_d(n)$, a functional annihilating $W_d(n)$ cannot
take value one on it. Conversely its class in the finite-dimensional
quotient $P_{\le\max\{d,D_G\}}(n)/W_d(n)$ is nonzero whenever
it is not in $W_d(n)$. Extend that class to a basis, assign it value
one and assign the other basis vectors zero. Composing with the
quotient map gives the required functional. If the target degree is
greater than $d$, it automatically lies outside $W_d(n)$; the enlarged
ambient space also covers this degree-floor obstruction. Failure of
all-target derivability is failure for at least one target.
\end{proof}

A $K$-valued solution $a$ of $\Gamma_n$ which violates $G_b^{(n)}=0$ gives a
witness by $\lambda(f)=f(a)/G_b^{(n)}(a)$. For target one this is ordinary
evaluation, with its multiplicative property. A general witness need
not be multiplicative, positive or normalised on one. It records a
linear separation from all permitted bounded-degree consequences,
not necessarily a truth assignment to the original variables. If the solution is instead defined over an extension field $E/K$, the scaled evaluation initially takes values in $E$. Choose a $K$-linear map $\pi:E\to K$ with $\pi(1)=1$, by extending $1$ to a $K$-basis of $E$. Composing the scaled evaluation with $\pi$ gives a $K$-valued separating functional: it still vanishes on every consequence and takes value one on the target. This descent uses linearity, not multiplicativity of $\pi$.

\begin{example}[Semantic entailment need not be derivation]
The equation $x^2=0$ over a field forces $x=0$ at every solution.
Nevertheless, neither NS nor PC derives $x$ from this axiom: every
consequence belongs to the ideal $(x^2)$. For $d\ge2$, the functional
on $P_{\le d}$ taking value one on $x$ and zero on all other monomials
annihilates the degree-$d$ consequence space. This is a target witness
although no actual solution violates $x=0$. The example explains why
our minimum degree is a syntactic ideal-membership resource.
\end{example}

\begin{example}[Linear consistency without truth-value evaluation]
Over $\mathbb F_3$, take $x^2-x=0$, target one and degree budget two.
The PC consequence space is $K(x^2-x)$: no nonzero vector in that span
has degree at most one, so no multiplication is usable. The functional
\[
 \lambda(1)=1,\qquad\lambda(x)=\lambda(x^2)=2
\]
is a witness but is not evaluation: $\lambda(x^2)=2$ whereas
$\lambda(x)^2=1$. Boolean equations impose the intended truth-value
semantics on actual solutions, without making every linear witness an
actual evaluation.
\end{example}

Annihilating only axiom multiples gives an NS witness. For PC the
functional must annihilate the full closure, including polynomials
obtained by cancellation followed by new multiplication. This is the
dual form of the distinction between $I_d$ and $C_d$.

\section{Four patterns of bounded-degree algebraic information}
Let $h(n)$ be a fixed-orbit or all-target degree profile in the chosen
system. By Proposition~\ref{tx:prop:dual-witness}, a degree-$d$ witness
exists precisely when $h(n)>d$.
Recall that $E(h)$ consists of the finite values occurring at arbitrarily
large universe sizes. The four alternatives classify eventual boundedness and
whether $E(h)$ is empty, finite nonempty or infinite; they do not classify
the sets of roots of $\Gamma_n$.

\begin{proposition}[A dual reading of the classification]
\label{tx:prop:semantic-types}
The four types have the following mutually exclusive interpretations.
\begin{enumerate}[label=\textup{T\arabic*.}]
\item Some fixed $D$ suffices to exclude degree-$D$ witnesses at every
      sufficiently large size.
\item For every fixed $d$, witnesses exist at every sufficiently large
      size.
\item There are a fixed $D$, a $p$-power period and a nonempty proper
      easy residue set. No degree-$D$ witness exists on easy residues
      at large sizes. On hard residues, witnesses exist at every fixed
      degree for all sufficiently large sizes in those residues.
\item Infinitely many distinct finite degrees $e$ recur, each at
      arbitrarily large sizes. At an occurrence with $e\ge1$, a degree-$(e-1)$ witness
      exists while a degree-$e$ witness does not. There is also a fixed
      degree at which witnesses fail at arbitrarily large sizes. No
      fixed periodic partition separates a uniformly bounded-degree
      part from a part with positive linear hardness.
\end{enumerate}
\end{proposition}
\begin{proof}
T1 follows directly from $h(n)>d$ being equivalent to witness
existence. In T2 the linear bound eventually exceeds any fixed $d$;
conversely eventual witness existence at every fixed degree excludes
every recurring finite value. For T3 use the easy/hard residue split
of Theorem~\ref{tx:thm:classification}. Conversely the displayed split
makes the recurring values a nonempty subset of $\{0,\ldots,D\}$
and leaves the profile unbounded. In T4 use each recurring $e$ for
the two adjacent-degree assertions, and use the least recurring value
for infinitely many failures at a fixed degree. The impossibility of
a fixed periodic bounded/linear split is
Proposition~\ref{tx:prop:no-periodic-split}. Infinitely many recurring
finite values characterise T4, completing the converse.
\end{proof}

For each fixed $d$, the existence of a witness is eventually determined
by a residue modulo a $p$-power. T3 needs only one finite degree window
to distinguish all easy residues. In T4, larger windows keep revealing
new recurring finite levels. The linear onset bounds the last possible
isolated low degree, while the short period locates finite linear peaks.
For indexed targets, the separating index may vary with size within
the fixed requested system; this is not a change to the request itself.

The \emph{finite-model spectrum} of a sentence is the set of positive
integers $n$ for which it has a model with $n$ elements. It records
existence of solutions, without recording a polynomial proof budget.

The refutation examples in Section~\ref{tx:sec:examples} all have no
finite nonempty models yet exhibit all four types. Their identical empty
finite-model spectra cannot distinguish the cases. The fresh-target
construction in Section~\ref{tx:sec:target-examples} also shows that a
nonconstant consequence can retain these algebraic distinctions. The
classification is a property of uniform algebraic reasoning about the
chosen equations and targets, rather than a satisfiability classification.

\section{Optional background: the Euler criterion}
\label{tx:sec:euler-background}
This section explains the separate sufficient criterion recalled at the
end of Section~11 of the paper. It is not another description of the
four types. Readers interested in the topological classification can
skip it without losing any step of the argument.

Kraj\'i\v cek's Euler structures~\cite{KrajEuler} equip the permitted
parameter-definable sets with a strong counting characteristic $\chi$.
It assigns one to a point, is additive on disjoint unions,
multiplicative on products and invariant under definable bijections.
It also satisfies the fibre rule $\chi(X)=r\chi(Y)$ for a definable
map $X\to Y$ whose fibres all have characteristic $r$. On finite sets
this extends ordinary counting; the permitted definable sets may be
infinite. The characteristic on definable sets is distinct from the
field characteristic $p$.

For the degree criterion the values lie in $\mathbb Z/p^\nu\mathbb Z$,
with precision $\nu$ depending on the degree budget and indexing bounds.
Reduction modulo $p$ interprets a definable $\mathbb F_p$-valued sum as
\[
 \sum_{x\in X}^{\chi}a(x)
       =\sum_{c\in\mathbb F_p}c\,\overline\chi(a^{-1}(c)).
\]
Under the full definability and index-sort hypotheses of
\cite[Theorem~6.1]{KrajEuler}, including the characteristic on the
required index sorts, a definable solution violating the target $g=0$
excludes derivations within the specified degree budget at sufficiently
large sizes in the residue class $\chi(M)\bmod p^\nu$.
The refutation case is $g=1$. A merely mod-$p$ counting rule or an
arbitrary infinite model does not supply these hypotheses. The
Grothendieck-ring viewpoint~\cite{KrajScanlon} organises the additive
and multiplicative identities; the strong fibre rule remains additional
structure.

This is a \emph{sufficient obstruction} to bounded-degree derivations
when the required structure exists. We do not prove that every hard
uniform family admits such a structure, nor identify T1--T4 with four
classes of Euler structures. A general converse remains a separate
question. The hardness-spectrum classification above needs no Euler
structure, no Boolean encoding and no model-existence assumption.

\section{What the model-theoretic language does not replace}
The finite-dimensional target theorem requires neither an infinite model
nor an Euler structure. Such structures can explain particular residue
obstructions, but the proof of the four types works directly with finite
consequence spaces and their actual maps. Boolean equations are essential
for the intended truth-value interpretation of the first-order encoding,
not for the general support-bounded transfer theorem.

A useful test of understanding is to say which statement is being proved:
that a target vanishes at all roots, that its polynomial lies in an ideal,
or that it has a derivation within one particular degree budget. These
are three different questions. The book's profile concerns the last.
